% GENERATED FILE. Edit canonical lessons, then run npm run generate:books.
% canonical-source-hash: fnv1a64:7d54049cd1c7a77f
% canonical-lessons: ZH-C395-shixian, ZH-C395-qude, ZH-W395-qu-take, ZH-C395-chengguo, ZH-C395-wancheng, ZH-C395-tiaojian

\chapter{Getting Results}
\label{ch:zh-getting-results}
\hlchaptermodality{395}

\begin{chapteropening}
\textbf{By the end of this chapter:} \emph{I can talk about dreams coming true, achieving something, results, finishing a job and a condition.}

Complete the last lesson of chapter 395: I can talk about dreams coming true, achieving something, results, finishing a job and a condition.
\end{chapteropening}

\section[shíxiàn]{shíxiàn — to come true; to make real}

\label{lesson:ZH-C395-shixian}



\begin{quote}
Before the new one: say the Mandarin for to be wasted (of effort), then the Mandarin for staying the same.
\end{quote}

\subsection*{You'll want to know: shíxiàn}
\emph{shíxiàn} — \textquotedblleft{}to come true; to make real\textquotedblright{}.

\begin{quote}
\emph{yàoshi yǒu qián wǒ de mèngxiǎng jiù néng shíxiàn} — If I had the money, my dream could come true.
\end{quote}

Its characters come later; for now it is a word for the ear.

\begin{cousinweb}
\emph{shíxiàn} is built from \emph{shí} 'real' and \zh{现} 'to appear', both inherited from Middle Chinese, \emph{zyit} and \emph{henH} (Guangyun, in Baxter's transcription), which glosses the first 'full' and 'true' and the second 'to show'.
\end{cousinweb}

\begin{grammarlens}[title={What you've built}]
One more word for saying what would happen.
\end{grammarlens}

\subsection*{Your turn}
\begin{itemize}
\raggedright
  \item \emph{Say it:} \emph{shíxiàn}
  \item \emph{Say it:} \emph{shíxiàn}, once more
  \item \emph{Say it:} \emph{yàoshi yǒu qián wǒ de mèngxiǎng jiù néng shíxiàn} — If I had the money, my dream could come true.
  \item \emph{Recall:} say the Mandarin for to be wasted (of effort), then the Mandarin for staying the same, then say \emph{shíxiàn} again
\end{itemize}

\begin{tcolorbox}[breakable,colback=teal!4,colframe=teal!35!black,title={Before you move on}]
What does shíxiàn mean? (\textquotedblleft{}To come true; to make real\textquotedblright{}.) Say it once more.
\end{tcolorbox}
\section[\texorpdfstring{\zh{取得} (qǔdé)}{取得 (qǔdé)}]{\zh{取得} (qǔdé) — to achieve; to obtain (through effort)}

\label{lesson:ZH-C395-qude}



\begin{quote}
Before the new one: say the Mandarin for staying the same, then the Mandarin for to come true.
\end{quote}

\subsection*{You'll want to know: \zh{取得}}
\textbf{\zh{取得}} — \emph{qǔdé} — \textquotedblleft{}to achieve; to obtain (through effort)\textquotedblright{}.

\begin{quote}
\emph{tā qǔdé le hěn hǎo de chéngjì} — She got very good marks.
\end{quote}

The new character, in pieces: \zh{取} — eight strokes in two pieces: on the left the ear piece, six strokes; on the right \zh{又}, as at the foot of \zh{变}. The next lesson puts it in your hand.

\begin{cousinweb}
\textbf{\zh{取得}} (\emph{qǔdé}) is built from \zh{取} 'to take' and \zh{得} 'to get', both inherited from Middle Chinese, \emph{tshjuX} and \emph{tok} (Guangyun, in Baxter's transcription), which glosses the first 'to gather in' and 'to receive'.
\end{cousinweb}

\begin{grammarlens}[title={What you've built}]
One more word for saying what would happen.
\end{grammarlens}

\subsection*{Your turn}
\begin{itemize}
\raggedright
  \item \emph{Say it:} \emph{qǔdé}
  \item \emph{Say it:} \emph{qǔdé}, once more
  \item \emph{Say it:} \emph{tā qǔdé le hěn hǎo de chéngjì} — She got very good marks.
  \item \emph{Recall:} say the Mandarin for staying the same, then the Mandarin for to come true, then say \emph{qǔdé} again
\end{itemize}

\begin{tcolorbox}[breakable,colback=teal!4,colframe=teal!35!black,title={Before you move on}]
What does \zh{取得} mean? (\textquotedblleft{}To achieve; to obtain (through effort)\textquotedblright{}.) Say it once more.
\end{tcolorbox}
\section[\texorpdfstring{\zh{取} (qǔ)}{取 (qǔ)}]{\zh{取} — the ear piece, then \zh{又}}

\label{lesson:ZH-W395-qu-take}



\begin{quote}
The word first, then the character inside it.

\begin{itemize}
\raggedright
  \item \emph{Recall:} say \emph{qǔdé} — \textbf{R1}, one lesson back, before it decays
\end{itemize}

\begin{quote}
\zh{变} \zh{又}
\end{quote}

\emph{Read it:} each one aloud, with nothing said first — \emph{biàn}, \emph{yòu}
\end{quote}

\subsection*{Script — the ear piece, then \zh{又}}
\hlblockfigure{\detokenize{figures/ZH-W395-qu-take-filmstrip.pdf}}{How \zh{取} is written, stroke by stroke}

\begin{quote}
\zh{取}
\end{quote}

\textbf{Two pieces, left before right.}

\textbf{The left piece, the ear piece, six strokes.} A bar across the top. A short vertical down on the left. A long vertical down on the right, past the foot. Two short bars between them. Then a stroke rising to the right at the foot.

\textbf{The right piece, \zh{又}, two strokes,} as at the foot of \zh{变}: a bar that turns and falls to the lower left, then a long stroke falling to the right across it.

6 + 2 = eight.

\subsection*{Your turn}
\begin{itemize}
\raggedright
  \item \emph{Write it:} \zh{取}, all eight, with the model in front of you
  \item \emph{Cover:} write it again from memory, and count the pieces, not the strokes
\end{itemize}

\begin{tcolorbox}[breakable,colback=teal!4,colframe=teal!35!black,title={Before you move on}]
How many strokes in \zh{取}? (\textbf{Eight}.) Which vertical is longer? (\textbf{The right one, past the foot}.) What is the right piece? (\textbf{\zh{又}, as at the foot of \zh{变}}.)

Source: \href{https://raw.githubusercontent.com/chanind/hanzi-writer-data/68d10a4b21150cae5e1ebbd223eed289cf32d90c/data/\%E5\%8F\%96.json}{Hanzi Writer Data \zh{取}, pinned stroke paths and medians}.
\end{tcolorbox}
\section[\texorpdfstring{\zh{成果} (chéngguǒ)}{成果 (chéngguǒ)}]{\zh{成果} (chéngguǒ) — fruits of work; achievements}

\label{lesson:ZH-C395-chengguo}



\begin{quote}
Before the new one: say the Mandarin for to come true, then the Mandarin for to obtain (through effort).
\end{quote}

\subsection*{You'll want to know: \zh{成果}}
\textbf{\zh{成果}} — \emph{chéngguǒ} — \textquotedblleft{}fruits of work; achievements\textquotedblright{}.

\begin{quote}
\textbf{\zh{我们今年取得了很好的成果}} \emph{wǒmen jīnnián qǔdé le hěn hǎo de chéngguǒ} — We got good results this year.
\end{quote}

\begin{cousinweb}
\textbf{\zh{成果}} (\emph{chéngguǒ}) is built from \zh{成} 'to complete' and \zh{果} 'fruit', both inherited from Middle Chinese, \emph{dzyeng} and \emph{kwaX} (Guangyun, in Baxter's transcription), which glosses the second 'the fruit of a tree' and 'resolute'.
\end{cousinweb}

\begin{grammarlens}[title={What you've built}]
One more word for saying what would happen.
\end{grammarlens}

\subsection*{Your turn}
\begin{itemize}
\raggedright
  \item \emph{Say it:} \emph{chéngguǒ}
  \item \emph{Say it:} \emph{chéngguǒ}, once more
  \item \emph{Say it:} \emph{wǒmen jīnnián qǔdé le hěn hǎo de chéngguǒ} — We got good results this year.
  \item \emph{Recall:} say the Mandarin for to come true, then the Mandarin for to obtain (through effort), then say \emph{chéngguǒ} again
\end{itemize}

\begin{tcolorbox}[breakable,colback=teal!4,colframe=teal!35!black,title={Before you move on}]
What does \zh{成果} mean? (\textquotedblleft{}Fruits of work; achievements\textquotedblright{}.) Say it once more.
\end{tcolorbox}
\section[wánchéng]{wánchéng — to complete; to get (a job) done}

\label{lesson:ZH-C395-wancheng}



\begin{quote}
Before the new one: say the Mandarin for to obtain (through effort), then the Mandarin for fruits of work.
\end{quote}

\subsection*{You'll want to know: wánchéng}
\emph{wánchéng} — \textquotedblleft{}to complete; to get (a job) done\textquotedblright{}.

\begin{quote}
\emph{rúguǒ jīntiān néng wánchéng míngtiān jiù kěyǐ xiūxi} — If we can finish today, we can rest tomorrow.
\end{quote}

Its characters come later; for now it is a word for the ear.

\begin{cousinweb}
\emph{wánchéng} is built from \emph{wán} 'whole' and \zh{成} 'to complete', both inherited from Middle Chinese, \emph{hwan} and \emph{dzyeng} (Guangyun, in Baxter's transcription), which glosses the first 'complete'.
\end{cousinweb}

\begin{grammarlens}[title={What you've built}]
One more word for saying what would happen.
\end{grammarlens}

\subsection*{Your turn}
\begin{itemize}
\raggedright
  \item \emph{Say it:} \emph{wánchéng}
  \item \emph{Say it:} \emph{wánchéng}, once more
  \item \emph{Say it:} \emph{rúguǒ jīntiān néng wánchéng míngtiān jiù kěyǐ xiūxi} — If we can finish today, we can rest tomorrow.
  \item \emph{Recall:} say the Mandarin for to obtain (through effort), then the Mandarin for fruits of work, then say \emph{wánchéng} again
\end{itemize}

\begin{tcolorbox}[breakable,colback=teal!4,colframe=teal!35!black,title={Before you move on}]
What does wánchéng mean? (\textquotedblleft{}To complete; to get (a job) done\textquotedblright{}.) Say it once more.
\end{tcolorbox}
\section[tiáojiàn]{tiáojiàn — condition (that must be met)}

\label{lesson:ZH-C395-tiaojian}



\begin{quote}
Before the new one: say the Mandarin for fruits of work, then the Mandarin for to get (a job) done.
\end{quote}

\subsection*{You'll want to know: tiáojiàn}
\emph{tiáojiàn} — \textquotedblleft{}condition (that must be met)\textquotedblright{}.

\begin{quote}
\emph{wǒ kěyǐ qù kěshì yǒu yí ge tiáojiàn} — I can go, but there's one condition.
\end{quote}

Its characters come later; for now it is a word for the ear.

\begin{cousinweb}
\emph{tiáojiàn} is built from \emph{tiáo} 'a twig' and \emph{jiàn} 'a portion', both inherited from Middle Chinese, \emph{dew} and \emph{gjenX} (Guangyun, in Baxter's transcription), which glosses the first 'a small branch' and the second 'to divide into items'.
\end{cousinweb}

\begin{grammarlens}[title={What you've built}]
One more word for saying what would happen.
\end{grammarlens}

\subsection*{Your turn}
\begin{itemize}
\raggedright
  \item \emph{Say it:} \emph{tiáojiàn}
  \item \emph{Say it:} \emph{tiáojiàn}, once more
  \item \emph{Say it:} \emph{wǒ kěyǐ qù kěshì yǒu yí ge tiáojiàn} — I can go, but there's one condition.
  \item \emph{Recall:} say the Mandarin for fruits of work, then the Mandarin for to get (a job) done, then say \emph{tiáojiàn} again
\end{itemize}

\begin{tcolorbox}[breakable,colback=teal!4,colframe=teal!35!black,title={Before you move on}]
What does tiáojiàn mean? (\textquotedblleft{}Condition (that must be met)\textquotedblright{}.) Say it once more.
\end{tcolorbox}
